\begin{abstract}
Recent advances in brain decoding have shown the potential of vision transformers (ViTs) to learn mappings from neural representations to visual features. Here we present CortexDecoder, a brain decoding framework that maps fMRI flatmaps to the images viewed by subjects. Flatmaps are neural responses projected onto the cortical surface and then unfolded into 2D images. We demonstrate that (1) flatmaps preserve sufficient signal to achieve high performance in brain decoding, whereas raw fMRI data is noisy and high-dimensional, which makes decoding difficult; and (2) reshaping raw fMRI data into flatmaps allows for clean patch-based tokenization. This lets us harness state-of-the-art ViTs to learn meaningful neural representations while preserving spatial encoding. To do this, we pass flatmap patches into a pretrained CortexMAE encoder to obtain brain tokens. A lightweight cross-attention bridge then learns to map these brain tokens to image features. The bridge is split into two modules: a semantic module that predicts CLIP embeddings and a low-level module that generates a low-level image reconstruction. Finally, a fine-tuned SDXL (Stable Diffusion XL) model takes this low-level reconstruction as its initial image and, conditioned on the predicted CLIP embedding, produces the final reconstruction. In evaluations on two NSD subjects, CortexDecoder achieves performance comparable to state-of-the-art brain decoding methods, demonstrating the efficacy of flatmap-based neural representations for visual reconstruction.
\end{abstract}
